% 原书物理页113--118；逐页图像自校，尚未独立复核.
\section{对称多项式}\label{sec:galois-2-9}
现在考虑 $n$ 个不定元的有理函数域 $K=k(a_1,\ldots,a_n)$，以及多项式 $f=X^{n-1}+a_1X^{n-1}+\cdots+a_n\in K[X]$. 我们将证明此多项式在 $K$ 上可分，且其 Galois 群为 $G_f=\mathfrak S_n$.\footnote{译注：原书此处首项印为 $X^{n-1}$；本节末尾讨论的首项则是 $X^n$. 此处保留原式.}
\subsection{对称群的作用}\label{sec:galois-2-9-1}
令 $\mathfrak S_n$ 按如下公式作用于 $\mathbb N^n$：
\begin{align*}
\sigma\cdot\nu=((\sigma\cdot\nu)_1,\ldots,(\sigma\cdot\nu)_n),\qquad (\sigma\cdot\nu)_i=\nu_{\sigma^{-1}(i)},\quad \nu=(\nu_1,\ldots,\nu_n).
\end{align*}
于是 $(\sigma\sigma')\cdot\nu=\sigma\cdot(\sigma'\cdot\nu)$，这说明我们定义了 $\mathfrak S_n$ 在 $\mathbb N^n$ 上的左作用.
由幺半群代数 $\mathbb Z[\mathbb N^n]$ 的泛性质，此作用延拓为 $\mathfrak S_n$ 通过环自同构在 $\mathbb Z[\mathbb N^n]$ 上的作用. 具体地，对于 $f=\sum_{\nu\in\mathbb N^n}a_\nu X^\nu$，有
\begin{align*}
\sigma(f)=\sum_{\nu\in\mathbb N^n}a_\nu X^{\sigma\cdot\nu}=\sum_{\nu\in\mathbb N^n}a_{\sigma^{-1}\cdot\nu}X^\nu.
\end{align*}
如同在 $\ldots\ldots$ 中一样，将 $\mathbb Z[\mathbb N^n]$ 与 $\mathbb Z[X_1,\ldots,X_n]$ 等同.\footnote{译注：原书此处的回指仅印作省略点，未给出目标.} 因而 $X^\nu=X_1^{\nu_1}\cdots X_n^{\nu_n}$，且
\begin{align*}
X^{\sigma\cdot\nu}=X_1^{\nu_{\sigma^{-1}(1)}}\cdots X_n^{\nu_{\sigma^{-1}(n)}}=X_{\sigma(1)}^{\nu_1}\cdots X_{\sigma(n)}^{\nu_n}.
\end{align*}
由此 $\sigma(X_i)=X_{\sigma(i)}$ 对所有 $i$ 成立. 换言之，环 $\mathbb Z[X_1,\ldots,X_n]$ 的自同构 $f\mapsto\sigma(f)$ 是满足 $\sigma(X_i):=X_{\sigma(i)}$ 的唯一自同构.
\subsection{初等对称多项式}\label{sec:galois-2-9-2}
对于 $j=1,\ldots,n$，置
\begin{align*}
\Sigma_j:=\sum_{1\leq i_1<\cdots<i_j\leq n}X_{i_1}\cdots X_{i_j}=\sum_{\substack{I\subset\{1,\ldots,n\}\\|I|=j}}X_I,\qquad X_I:=\prod_{i\in I}X_i.
\end{align*}
例如 $\Sigma_1=X_1+\cdots+X_n$，$\Sigma_n=X_1\cdots X_n$. 注意 $\mathfrak S_n$ 按 $(\sigma,I)\mapsto\sigma(I)$ 作用于 $\{1,\ldots,n\}$ 的子集，且显然保持基数. 又有 $\sigma(X_I)=X_{\sigma(I)}$，故对每个置换 $\sigma\in\mathfrak S_n$，有 $\sigma(\Sigma_n)=\Sigma_n$. 换言之，
\begin{align*}
\Sigma_1,\ldots,\Sigma_n\in\mathbb Z[X_1,\ldots,X_n]^{\mathfrak S_n}.
\end{align*}
我们称它们为``对称多项式''. 这些多项式编码了多项式根与系数之间的关系.
\begin{lemma*}在 $\mathbb Z[X_1,\ldots,X_n][T]$ 中有
\begin{align*}
(T-X_1)\cdots(T-X_n)=T^n-\Sigma_1T^{n-1}+\Sigma_2T^{n-2}+\cdots+(-1)^n\Sigma_n.
\end{align*}
\end{lemma*}
\begin{proof}留给读者对 $n$ 归纳证明.\end{proof}
给定交换环 $A$ 中的 $n$ 元组 $\alpha_1,\ldots,\alpha_n$，以将 $T$ 映为 $X$、将 $\alpha_i$ 映为 $X_i$ 的同态
\begin{align*}\mathbb Z[X_1,\ldots,X_n][T]\to A[X]\end{align*}
特化引理中的恒等式，\footnote{译注：原书此处将 $X_i$ 与 $\alpha_i$ 的方向写反；该同态应将 $X_i$ 映为 $\alpha_i$.} 在 $A[X]$ 中得到
\begin{align*}
(X-\alpha_1)\cdots(X-\alpha_n)=X^n-\Sigma_1(\alpha_1,\ldots,\alpha_n)X^{n-1}+\cdots+(-1)^n\Sigma_n(\alpha_1,\ldots,\alpha_n).
\end{align*}
下面这个出色的定理说明了\newterm{galois-elementary-symmetric-polynomial}{初等对称多项式}{elementary symmetric polynomial}这一名称的由来.
\subsection{对称多项式定理}\label{sec:galois-2-9-3}
\begin{theorem*}$\Sigma_i$ 在 $\mathbb Z[X_1,\ldots,X_n]$ 中代数无关，并生成子环 $\mathbb Z[X_1,\ldots,X_n]^{\mathfrak S_n}$. 换言之，将 $Y_i$ 映为 $\Sigma_i$ 的唯一环同态 $\mathbb Z[Y_1,\ldots,Y_n]\to\mathbb Z[X_1,\ldots,X_n]$ 是单射，其像为 $\mathbb Z[X_1,\ldots,X_n]^{\mathfrak S_n}$，略不严格地写作
\begin{align*}
\mathbb Z[X_1,\ldots,X_n]^{\mathfrak S_n}=\mathbb Z[\Sigma_1,\ldots,\Sigma_n].
\end{align*}\end{theorem*}
\begin{proof}
对于 $\nu\in\mathbb N^n$，置 $\Sigma^\nu:=\Sigma_1^{\nu_1}\cdots\Sigma_n^{\nu_n}$. 我们须证明 $(\Sigma^\nu)_{\nu\in\mathbb N^n}$ 是 $\mathbb Z[X_1,\ldots,X_n]^{\mathfrak S_n}$ 的一组 $\mathbb Z$-基.
为此，先给出这个环的一组方便的 $\mathbb Z$-基，再用它表示 $\Sigma^\nu$. 设 $f=\sum_{\nu\in\mathbb N^n}a_\nu X^\nu$ 在 $\mathfrak S_n$ 下不变. 则对每个 $\sigma\in\mathfrak S_n$，有 $a_\nu=a_{\sigma\cdot\nu}$，因而可以唯一地写成
\begin{align*}
f=\sum_{\bar\nu\in\mathbb N^n/\mathfrak S_n}a_{\bar\nu}S^{\bar\nu},\qquad S^{\bar\nu}:=\sum_{\nu\in\bar\nu}X^\nu.
\end{align*}
故 $(S^{\bar\nu})_{\bar\nu\in\mathbb N^n/\mathfrak S_n}$ 是所求的一组 $\mathbb Z$-基. 但用商集 $\mathbb N^n/\mathfrak S_n$ 参数化这组基，不便于表示 $\Sigma^\nu$.
\begin{definition*}若 $\nu_1\geq\nu_2\geq\cdots\geq\nu_n$，则称 $\nu\in\mathbb N^n$ 是优势的. 以 $\Lambda\subset\mathbb N^n$ 表示优势 $n$ 元组的集合.\end{definition*}
每个 $\mathfrak S_n$-轨道显然恰含一个优势 $n$ 元组. 因而 $\Lambda$ 是 $\mathbb N^n$ 中各 $\mathfrak S_n$-轨道的一组代表元. 对于 $\lambda\in\Lambda$，简记
\begin{align*}
S^\lambda:=S^{\bar\lambda}=\sum_{\nu\in\mathfrak S_n\cdot\lambda}X^\nu.
\end{align*}
\begin{example*}对 $j=1,\ldots,n$，令 $\mu^j=(1,\ldots,1,0,\ldots,0)\in\Lambda$，其中有 $j$ 项等于 $1$. 则 $\mu^j\in\Lambda$ 且 $S^{\mu^j}=\Sigma_j$.\end{example*}
若 $\lambda'\in\Lambda$，则可写成
\begin{align*}
S^\lambda S^{\lambda'}=\sum_{\mu\in\Lambda}c_{\lambda,\lambda';\mu}S^\mu.\tag{$*$}\label{eq:galois-symmetric-product}
\end{align*}
其中系数 $c_{\lambda,\lambda';\mu}\in\mathbb Z$ 唯一确定. 为研究这些系数，注意 $\Lambda$ 对加法封闭，并在 $\mathbb N^n$ 上引入如下次序.
\begin{definition*}对 $\nu\in\mathbb N^n$，置 $|\nu|:=\sum_{i=1}^n\nu_i$，$\widetilde\nu:=(|\nu|,\nu_1,\ldots,\nu_n)\in\mathbb N^{n+1}$. 在 $\mathbb N^n$ 上定义次序关系
\begin{align*}
\nu\preccurlyeq\nu'\quad\Longleftrightarrow\quad\widetilde\nu\leq\widetilde\nu',
\end{align*}
其中 $\leq$ 是 $\mathbb N^{n+1}$ 上的字典序.\end{definition*}
次序 $\preccurlyeq$ 是全序且与加法相容：$\nu\preccurlyeq\nu'$ 且 $\mu\preccurlyeq\mu'$ 蕴含 $\nu+\mu\preccurlyeq\nu'+\mu'$.
\begin{lemma*}在展开式~\eqref{eq:galois-symmetric-product} 中，$c_{\lambda,\lambda';\mu}\ne0$ 蕴含 $\mu\preccurlyeq\lambda+\lambda'$，且 $c_{\lambda,\lambda';\lambda+\lambda'}=1$. 换言之，
\begin{align*}
S^\lambda S^{\lambda'}\in S^{\lambda+\lambda'}+\sum_{\mu\prec\lambda+\lambda'}\mathbb Z S^\mu.
\end{align*}\end{lemma*}
\begin{proof}对 $\nu\in\mathbb N^n$，记 $A_{\prec\nu}:=\sum_{\nu'\prec\nu}\mathbb ZX^{\nu'}$. 因为 $\preccurlyeq$ 与加法相容，有 $A_{\prec\nu}A_{\preccurlyeq\nu'}\subset A_{\prec\nu+\nu'}$. 又若 $\lambda\in\Lambda$，则对每个 $\sigma\in\mathfrak S_n\setminus\{\mathrm{id}\}$，有 $\sigma\cdot\lambda\prec\lambda$.\footnote{译注：此断言在 $\lambda$ 有重复分量时不成立；非恒等置换可能固定 $\lambda$. 原书此处保留，后面的轨道和按不同轨道元素取和.} 所以 $S^\lambda\in X^\lambda+A_{\prec\lambda}$，且
\begin{align*}
S^\lambda S^{\lambda'}\in X^{\lambda+\lambda'}+A_{\prec\lambda+\lambda'}.
\end{align*}
因此 $S^\lambda S^{\lambda'}=S^{\lambda+\lambda'}+P$，其中 $P\in(A_{\prec\lambda+\lambda'})^{\mathfrak S_n}$. 写 $P=\sum_{\mu\in\Lambda}c_\mu S^\mu$. 因 $c_\mu$ 也是将 $P$ 在 $X^\nu$ 基下展开时 $X^\mu$ 的系数，故 $c_\mu\ne0$ 蕴含 $\mu\prec\lambda+\lambda'$，如所欲证.\end{proof}
现在用 $S^\lambda$ 基表示 $\Sigma^\nu$. 有
\begin{align*}
\Sigma^\nu=\Sigma_1^{\nu_1}\cdots\Sigma_n^{\nu_n}=(S^{\mu^1})^{\nu_1}\cdots(S^{\mu^n})^{\nu_n}.
\end{align*}
置 $\lambda_\nu:=\nu_1\mu^1+\cdots+\nu_n\mu^n\in\Lambda$. 引理给出
\begin{align*}
\Sigma^\nu\in S^{\lambda_\nu}+\sum_{\lambda\prec\lambda_\nu}\mathbb Z S^\lambda.
\end{align*}
而映射 $\mathbb N^n\to\Lambda,\ \nu\mapsto\lambda_\nu$ 是双射，其逆为 $\lambda\mapsto(\lambda_1-\lambda_2,\ldots,\lambda_{n-1}-\lambda_n,\lambda_n)$. 因此在 $\mathbb N^n$ 上引入新的次序关系（也是与加法相容的全序）：$\nu\trianglelefteq\nu'\Longleftrightarrow\lambda_\nu\preccurlyeq\lambda_{\nu'}$，并置 $S_\nu:=S^{\lambda_\nu}$. 则 $(S_\nu)_{\nu\in\mathbb N^n}$ 是 $\mathbb Z[X_1,\ldots,X_n]^{\mathfrak S_n}$ 的一组 $\mathbb Z$-基，且
\begin{align*}
\forall\nu\in\mathbb N^n,\qquad\Sigma^\nu\in S_\nu+\sum_{\nu'\triangleleft\nu}\mathbb Z S_{\nu'}.
\end{align*}
由此 $(\Sigma^\nu)_{\nu\in\mathbb N^n}$ 在 $\mathbb Z$ 上线性无关. 事实上，若 $\sum_\nu a_\nu\Sigma^\nu=0$ 是线性相关关系，且 $\{\nu:a_\nu\ne0\}$ 非空，则此集合有最大元 $\nu_0$，于是得到关系 $0\in a_{\nu_0}S_{\nu_0}+\sum_{\nu\triangleleft\nu_0}S_\nu$，与 $S_\nu$ 的线性无关性矛盾.\footnote{译注：原书末式求和项未写系数，原样保留.}
还可推出 $(\Sigma^\nu)_{\nu\in\mathbb N^n}$ 在 $\mathbb Z$ 上生成整个环. 更确切地，例如对整数 $n(\nu)=|\{\nu'\triangleleft\nu\}|$ 归纳可得
\begin{align*}
\forall\nu\in\mathbb N^n,\qquad S_\nu\in\Sigma^\nu+\sum_{\nu'\triangleleft\nu}\mathbb Z\Sigma^{\nu'}.
\end{align*}
\end{proof}
\subsection{应用：多项式的判别式}\label{sec:galois-2-9-4}
由上述定理，存在唯一多项式 $\Delta\in\mathbb Z[\Sigma_1,\ldots,\Sigma_n]$ 满足
\begin{align*}
\prod_{i<j}(X_i-X_j)^2=(-1)^{n(n-1)/2}\prod_{i\ne j}(X_i-X_j)=\Delta(\Sigma_1,\ldots,\Sigma_n).
\end{align*}
因为左端显然是关于 $X_i$ 的对称多项式.
\begin{definition*}设 $A$ 是环，$f=X^n+a_1X^{n-1}+\cdots+a_{n-1}X+a_n\in A[X]$ 是首一多项式. 定义 $f$ 的\newterm{galois-discriminant}{判别式}{discriminant}为
\begin{align*}
\operatorname{Disc}(f):=\Delta(-a_1,\ldots,(-1)^na_n)\in A.
\end{align*}\end{definition*}
\begin{example*}令 $A=\mathbb Z[X_1,\ldots,X_n]$，并令 $f_{\mathrm{univ}}$ 为次数 $n$ 的``泛''分裂多项式
\begin{align*}
f_{\mathrm{univ}}:=(T-X_1)\cdots(T-X_n)\in\mathbb Z[X_1,\ldots,X_n][T].
\end{align*}
则 $\operatorname{disc}(f_{\mathrm{univ}})=\Delta$，因为 $f_{\mathrm{univ}}=T^n-\Sigma_1T^{n-1}+\cdots+(-1)^n\Sigma_n$. 计算还可得到
\begin{align*}
\operatorname{disc}(f_{\mathrm{univ}})=(-1)^{n(n-1)/2}\prod_{i=1}^nf'(X_i).
\end{align*}\end{example*}
\begin{proposition*}设 $k$ 是域，$f\in k[X]$ 首一. 对于使得 $f$ 在 $K[X]$ 中分裂为 $f=(X-\alpha_1)\cdots(X-\alpha_n)$ 的任意扩张 $K$，有
\begin{align*}
\operatorname{disc}(f)=\prod_{i<j}(\alpha_i-\alpha_j)^2=(-1)^{n(n-1)/2}\prod_{i=1}^nf'(\alpha_i).
\end{align*}
特别地，$f$ 可分当且仅当 $\operatorname{disc}(f)\ne0$.\end{proposition*}
\begin{proof}对泛例子施以将 $T$ 映为 $X$、$X_i$ 映为 $\alpha_i$ 的同态 $\mathbb Z[X_1,\ldots,X_n][T]\to K[X]$ 即得.\end{proof}
\begin{exercise*}证明 $\operatorname{disc}(X^n+aX+b)=(-1)^{n(n-1)/2}((1-n)^{n-1}a^n+n^nb^{n-1})$.\end{exercise*}
下面是判别式的另一个应用.
\begin{proposition*}设 $f\in k[X]$ 可分，将其 Galois 群 $G_f$ 看作其在分裂域 $K_f$ 中各根的置换群 $\mathfrak S_n$ 的子群. 若 $\operatorname{car}(k)\ne2$，则 $G_f\subset\mathfrak A_n\Longleftrightarrow\operatorname{disc}(f)\in(k^\times)^2$.\end{proposition*}
\begin{proof}设 $\alpha_1,\ldots,\alpha_n$ 是 $f$ 在 $K_f$ 中的根. 将此编号给出的嵌入 $G_f\hookrightarrow\mathfrak S_n$ 记作 $\tau\mapsto\sigma_\tau$. 于是对所有 $i$，有 $\tau(\alpha_i)=\alpha_{\sigma_\tau(i)}$.
置 $D:=\prod_{i<j}(\alpha_i-\alpha_j)\in K_f$. $\operatorname{disc}(f)$ 在 $K_f$ 中的两个平方根是 $D$ 与 $-D$. 故 $\operatorname{disc}(f)\in k^2\Longleftrightarrow D\in k$. 因 $k=K_f^{G_f}$，考察 $G_f$ 在 $D$ 上的作用. 以 $\varepsilon:\mathfrak S_n\to\{\pm1\}$ 记符号，对每个 $\tau\in G_f$，有
\begin{align*}
\tau(D)=\prod_{i<j}(\alpha_{\sigma_\tau(i)}-\alpha_{\sigma_\tau(j)})=\varepsilon(\sigma_\tau)\prod_{i<j}(\alpha_i-\alpha_j)=\varepsilon(\sigma_\tau)D.
\end{align*}
若 $\operatorname{car}(k)\ne2$，由此得到 $D\in K_f^{G_f}=k\Longleftrightarrow G_f\subset\mathfrak A_n$.\end{proof}
\subsection{应用：一般多项式的 Galois 群}\label{sec:galois-2-9-5}
先给出定理~\ref{sec:galois-2-9-3} 关于对称有理函数的一个推论. 顺便注意，$\mathfrak S_n$ 在 $k[X_1,\ldots,X_n]$ 上的作用唯一地延拓至 $k(X_1,\ldots,X_n)$.
\begin{corollary*}有 $k(X_1,\ldots,X_n)^{\mathfrak S_n}=k(\Sigma_1,\ldots,\Sigma_n)$.\end{corollary*}
\begin{proof}包含 $\supset$ 显然. 对另一包含，取 $\phi\in k(X_1,\ldots,X_n)^{\mathfrak S_n}$. 由于 $k[X_1,\ldots,X_n]$ 是唯一分解整环，可将 $\phi$ 唯一写成 $\phi=f/g$，其中 $(f,g)=1$，且 $f,g$ 首一.\footnote{译注：原书如此断言，但一般有理函数的互素分子、分母不能同时归一为首一；多元情形也须指定首项次序. 此处保留原论证，未补证明.} 对任意 $\sigma\in\mathfrak S_n$，有 $\phi=\sigma(\phi)=\sigma(f)/\sigma(g)$，且 $(\sigma(f),\sigma(g))=1$，$\sigma(f),\sigma(g)$ 首一. 由此 $f=\sigma(f)$，$g=\sigma(g)$. 所以 $f,g\in k[\Sigma_1,\ldots,\Sigma_n]$，最终 $\phi\in k(\Sigma_1,\ldots,\Sigma_n)$.\end{proof}
现在换个观点. 设 $k$ 是域，$K:=k(a_1,\ldots,a_n)$ 是 $n$ 个不定元 $a_1,\ldots,a_n$ 的有理函数域. 考察``一般''多项式 $f=X^n+a_1X^{n-1}+\cdots+a_n\in K[X]$，我们希望求其分裂域与 Galois 群 $G_f$. 为此，考虑 $n$ 个不定元 $\alpha_1,\ldots,\alpha_n$ 的有理函数域 $L:=k(\alpha_1,\ldots,\alpha_n)$. 定理~\ref{sec:galois-2-9-3} 说明 $\Sigma_i(\alpha_1,\ldots,\alpha_n)\in L$ 在 $k$ 上代数无关. 因而存在唯一嵌入
\begin{align*}
K\hookrightarrow L,\qquad a_i\mapsto(-1)^i\Sigma_i(\alpha_1,\ldots,\alpha_n).
\end{align*}
\begin{theorem*}多项式 $f$ 在 $K[X]$ 中可分且不可约. $L$ 是 $f$ 在 $K$ 上的一个分裂域，$\alpha_i$ 是其中的各根. Galois 群 $G_f$ 在 $\alpha_i$ 上的作用将 $G_f$ 等同于 $\mathfrak S_n$.\end{theorem*}
\begin{proof}由将 $X_i$ 映为 $\alpha_i$、$T$ 映为 $X$ 的环同态 $\mathbb Z[X_1,\ldots,X_n][T]\to L[X]$，引理~\ref{sec:galois-2-9-2} 中的分解给出 $L[X]$ 中的分解 $(X-\alpha_1)\cdots(X-\alpha_n)=X^n+a_1X^{n-1}+\cdots+a_n$. 因而 $\alpha_i$ 是 $f$ 在 $L$ 中的各根，所以 $f$ 可分且 $L$ 是其分裂域. $\operatorname{Gal}(L/K)$ 在 $\alpha_i$ 上的作用给出嵌入 $\operatorname{Gal}(L/K)\hookrightarrow\mathfrak S_n$. 而上述推论给出 $K=L^{\mathfrak S_n}$，推论~\ref{sec:galois-2-7-6} 蕴含 $\operatorname{Gal}(L/K)=\mathfrak S_n$. 特别地，$\operatorname{Gal}(L/K)$ 在 $f$ 的根上传递作用，故由引理~\ref{sec:galois-2-8-1}，$f$ 不可约.\end{proof}
\begin{remark*}这样的结果对于我们真正关心的、$a_i$ 为 $k$ 中元素的多项式有何含义？答案强烈依赖于 $k$. 例如，若 $k=\mathbb C$，将 $a_i$ 特化为 $\mathbb C$ 中元素得到的每个多项式 $f$ 都分裂，故其 Galois 群平凡！若 $k=\mathbb R$ 且 $n>2$，$f$ 的特化从不不可约，且其 Galois 群平凡或等于 $\mathbb Z/2\mathbb Z$. 若 $k=\mathbb F_p$，$f$ 的特化可以不可约，但其 Galois 群总是 Abel 群. 而当 $k=\mathbb Q$ 时，Hilbert 的一个结果断言，$f$ 有无穷多个特化，使特化后的多项式不可约且 Galois 群为 $\mathfrak S_n$！\end{remark*}

